\section*{Chapitre \ref{ch:m2:government-bonds} --- Obligations d'État}

\begin{solution}{exo:m2:government-bonds:1}
Quarante et un jours sur 184 se sont écoulés~: $A = 2.125 \times 41/184 = 0.473505$ pour
100.
\end{solution}

\begin{solution}{exo:m2:government-bonds:2}
$99 + 16.5/32 = 99.515625$. $101.296875 = 101 + 9.5/32$~: 101-09+.
\end{solution}

\begin{solution}{exo:m2:government-bonds:3}
À 4~\%~: 100 (\cref{prop:m2:government-bonds:priceyield}). À 5~\%~: avec $v =
1/1.025$, $2\sum_{k=1}^{6} v^k + 100 v^6 = 97.2459$.
\end{solution}

\begin{solution}{exo:m2:government-bonds:4}
Flux de 2, 2, 2 et 102 aux échéances de 1, 2, 3 et 4 semestres, $v = 1/1.02$~: \omterm{def:m2:government-bonds:duration}{duration de Macaulay}
1.9419 ans, \omterm{def:m2:government-bonds:duration}{duration modifiée} $1.9419/1.02 = 1.9039$, \omterm{def:m2:government-bonds:duration}{DV01} $100 \times
1.9039 \times 10^{-4} = 0.019039$ pour 100, 190.39~USD par million.
\end{solution}

\begin{solution}{exo:m2:government-bonds:5}
Duration seule~: $+100.870911 \times 7.974966 \times 0.01 = +8.0444$. Avec
\omterm{def:m2:government-bonds:duration}{convexité}~: $+8.0444 + \frac12 \times 100.870911 \times 75.824 \times 0.0001 =
+8.4268$. Réévaluation complète à 3.20~\%~: $+8.4403$. Pour la baisse, comme pour la hausse, la
estimation au second ordre reste à moins de 0.014 de la valeur exacte~; celle au premier
ordre s'écarte de 0.40.
\end{solution}

\begin{solution}{exo:m2:government-bonds:6}
Actual/actual~: 44 jours d'une période de 181 jours, $2.5 \times 44/181 = 0.607735$.
30/360~: du 15 février au 31 mars comptent $30 + 16 = 46$ jours, $5 \times 46/360
= 0.638889$. L'écart, de 0.031 pour 100 (312~USD par million), importe à
quiconque détient une obligation d'une convention, couverte ou financée par
un instrument de l'autre~: les modalités de chaque obligation précisent sa \omterm{def:m2:government-bonds:daycount}{convention de décompte des jours}, et un
système qui en suppose une pour toutes valorisera mal les autres.
\end{solution}

\begin{solution}{exo:m2:government-bonds:7}
Facteur d'actualisation à dix ans 0.657651, \omterm{def:m2:government-bonds:zero}{taux zéro-coupon} 4.2350~\%, contre un
\omterm{def:m2:government-bonds:zero}{taux au pair} de 4.21~\%. La courbe des \omterm{def:m2:government-bonds:zero}{taux au pair} croît de deux à dix ans, si bien que les
premiers \omterm{def:m2:government-bonds:coupon}{coupons} d'une obligation au \omterm{def:m2:government-bonds:coupon}{pair} à dix ans sont actualisés à des \omterm{def:m2:government-bonds:zero}{taux zéro-coupon} plus bas que
celui de son dernier paiement~; son rendement est une moyenne pondérée de ces taux, inférieure au
\omterm{def:m2:government-bonds:zero}{taux zéro-coupon} à dix ans.
\end{solution}

\begin{solution}{exo:m2:government-bonds:8}
Seule une obligation \omterm{def:m2:government-bonds:coupon}{zéro-coupon} a une duration égale à sa durée de vie. Une
obligation de 4.5~\% à trente ans au \omterm{def:m2:government-bonds:coupon}{pair} a une \omterm{def:m2:government-bonds:duration}{duration de Macaulay} de 16.74 ans et une
\omterm{def:m2:government-bonds:duration}{duration modifiée} de 16.37~; une hausse d'un point de base coûte $16.37 \times 0.0001 = 0.164\%$ de son prix,
1\,637~USD par million, soit environ la moitié de l'affirmation.
\end{solution}

\begin{solution}{pb:m2:government-bonds:1}
\textbf{1.} L'échéance, le 30 septembre, tombe en fin de mois, donc les dates de \omterm{def:m2:government-bonds:coupon}{coupon}
sont le 31 mars et le 30 septembre~: précédent le 31 mars 2026, suivant le 30 septembre 2026,
cinq jours après le règlement-livraison.
\textbf{2.} $A = 2 \times 178/183 = 1.945355$~; \omterm{def:m2:government-bonds:accrued}{prix pied de coupon} 100.224893, coté
100-07, arrondi au 64e le plus proche.
\textbf{3.} \omterm{def:m2:government-bonds:accrued}{Prix plein coupon} 102.170249~: 102\,170\,249~USD.
\textbf{4.} À cinq ans~: 451.43~USD par million~; à dix ans~: 804.44~USD par million.
\textbf{5.} 23.18 et 75.82.
\textbf{6.} $100 \times 451.43/804.44 = \text{USD}~56.12$ millions de l'obligation à dix ans
(rapport 0.5612).
\textbf{7.} La valeur en \omterm{def:m2:government-bonds:accrued}{prix plein coupon} de la jambe vendeuse est $56\,117\,663 \times 1.00870911 =
\text{USD}~56\,606\,398$~; net $102\,170\,249 - 56\,606\,398 = \text{USD}~45\,563\,851$
longue, à financer (\cref{ch:m2:repo-and-specials}).
\textbf{8.} $-\text{USD}~9.61$~: la position n'a pas de \omterm{def:m2:government-bonds:duration}{DV01}, seulement une trace de
\omterm{def:m2:government-bonds:duration}{convexité}.
\textbf{9.} $+25$~: $-\text{USD}~5\,934$. $-25$~: $-\text{USD}~6\,092$.
\textbf{10.} La position est courte en \omterm{def:m2:government-bonds:duration}{convexité}~: par unité de \omterm{def:m2:government-bonds:duration}{DV01}, l'obligation à dix ans
présente davantage de \omterm{def:m2:government-bonds:duration}{convexité} ($75.8/7.97$, contre $23.2/4.42$), et le desk est vendeur
de celle-ci. Quel que soit le sens du mouvement, le prix de la jambe vendeuse varie
davantage que ne le prédit sa duration, dans le sens défavorable.
\textbf{11.} $-\text{USD}~450\,252$.
\textbf{12.} $+\text{USD}~449\,295$.
\textbf{13.} Une hausse du rendement à dix ans par rapport à celui à cinq ans~: la
position est un 5s10s \emph{steepener}.
\textbf{14.} Son \omterm{def:m2:government-bonds:duration}{DV01} est nul par construction, de même que sa duration (\omterm{def:m2:government-bonds:duration}{DV01}
divisé par la valeur)~: les 45.6 millions de valeur nette n'entraînent aucun
risque de taux du premier ordre.
\textbf{15.} Non. Le temps raccourcit les deux titres, et le \omterm{def:m2:government-bonds:coupon}{coupon} de l'obligation à cinq
ans du 30 septembre quitte son \omterm{def:m2:government-bonds:accrued}{prix plein coupon}~; après 28 jours, le \omterm{def:m2:government-bonds:duration}{DV01} de la jambe longue est
44\,519~USD et celui de la jambe vendeuse 44\,863~USD~: position nette courte de 344~USD par point de base.
La couverture doit être rééquilibrée.
\textbf{16.} Un point de base fait bouger un titre à dix ans presque deux fois autant
qu'un titre à cinq ans~: à nominal ou à valeur égale, on laisserait une forte position
directionnelle. Le \omterm{def:m2:government-bonds:duration}{DV01} correspond à la variation de prix par unité de variation de rendement.
\textbf{17.} Risque de courbe (les deux rendements n'évoluent pas de la même façon), \omterm{def:m2:government-bonds:duration}{convexité}, et la
dérive des \omterm{def:m2:government-bonds:duration}{DV01}~; s'ajoute aussi le financement de chaque jambe en repo (\cref{ch:m2:repo-and-specials}).
\textbf{18.} Les cotations et les prix de transaction sont en 32es par convention~; le risque,
P\&L et ratios de couverture sont des quantités continues qu'il ne faut pas arrondir au
64e avant de les sommer.
\textbf{19.} Résultat retenu~: \emph{la couverture neutre en \omterm{def:m2:government-bonds:duration}{DV01}} vend 56.12~millions de dollars
de l'obligation à dix ans contre 100~millions de dollars de celle à cinq ans (rapport 0.5612)~;
un mouvement parallèle de 25 points de base, dans un sens ou dans l'autre, coûte environ 6\,000~USD
($-5\,934$ à la hausse, $-6\,092$ à la baisse), le prix d'être court en \omterm{def:m2:government-bonds:duration}{convexité}, tandis
qu'une pentification de 10 points de base rapporte 449\,295~USD.
\textbf{20.} Sur la forme de la courbe, et non sur son niveau.
\end{solution}

\begin{solution}{iq:m2:government-bonds:1}
Le \omterm{def:m2:government-bonds:accrued}{prix plein coupon} est ce que paie l'acheteur~: le \omterm{def:m2:government-bonds:accrued}{prix pied de coupon} augmenté du
\omterm{def:m2:government-bonds:accrued}{coupon couru} depuis le dernier \omterm{def:m2:government-bonds:coupon}{coupon}, qui revient au vendeur. Le \omterm{def:m2:government-bonds:accrued}{prix pied de coupon}
supprime la dent de scie du couru et du paiement des \omterm{def:m2:government-bonds:coupon}{coupons}, si bien que ses
variations ne traduisent que celles du rendement~; c'est pourquoi les écrans l'affichent et
pourquoi les rendements se calculent à partir de lui.
\iqlookfor{le \omterm{def:m2:government-bonds:accrued}{coupon couru} comme différence, et pourquoi un prix coté ne doit
pas sauter aux dates de \omterm{def:m2:government-bonds:coupon}{coupon}.}
\end{solution}

\begin{solution}{iq:m2:government-bonds:2}
Le prix est une somme de facteurs d'actualisation $(1 + y/f)^{-t}$, chacun convexe en $y$.
La \omterm{def:m2:government-bonds:duration}{convexité} signifie que l'obligation gagne davantage à une baisse du rendement qu'elle ne perd à une
hausse égale, ce qui est favorable au détenteur, toutes choses égales par ailleurs~; les marchés la valorisent,
si bien qu'une position longue en \omterm{def:m2:government-bonds:duration}{convexité} sacrifie en général du rendement (carry) en contrepartie.
\iqlookfor{la somme de termes convexes, et le compromis entre carry et \omterm{def:m2:government-bonds:duration}{convexité}.}
\end{solution}

\begin{solution}{iq:m2:government-bonds:3}
Le \omterm{def:m2:government-bonds:coupon}{zéro-coupon}~: sa duration vaut dix ans, celle de l'obligation à \omterm{def:m2:government-bonds:coupon}{coupons} est plus courte (environ
huit ans pour un rendement de 5~\%), car les \omterm{def:m2:government-bonds:coupon}{coupons} versés plus tôt ramènent le temps moyen vers l'avant.
Pour 100 de nominal, la comparaison des \omterm{def:m2:government-bonds:duration}{DV01} exige aussi les prix~: le \omterm{def:m2:government-bonds:coupon}{zéro-coupon} se traite
très en dessous du \omterm{def:m2:government-bonds:coupon}{pair} (environ 61 à 5~\%), si bien que son \omterm{def:m2:government-bonds:duration}{DV01} pour 100 de nominal ($61 \times 9.75 \times
10^{-4} \approx 0.059$) est plus petit que celui de l'obligation à \omterm{def:m2:government-bonds:coupon}{coupons} ($100 \times 7.8 \times
10^{-4} \approx 0.078$). La duration se rapporte à une unité de valeur, le \omterm{def:m2:government-bonds:duration}{DV01} à une unité de nominal.
\iqlookfor{distinguer la duration (relative) du \omterm{def:m2:government-bonds:duration}{DV01} (absolu).}
\end{solution}

\begin{solution}{iq:m2:government-bonds:4}
Vendre le dix ans en nominal $N_{10} = N_5 \times \mathrm{DV01}_5/\mathrm{DV01}_{10}$,
soit environ 0.55 à 0.6 du nominal du cinq ans aux rendements actuels. Restent~:
le risque de courbe (on est désormais en steepener), la \omterm{def:m2:government-bonds:duration}{convexité} (position vendeuse), la dérive du \omterm{def:m2:government-bonds:duration}{DV01} dans le temps
et avec les rendements, le financement par repo et le specialness de chaque jambe, et la
fourchette lors du rééquilibrage.
\iqlookfor{rapport des \omterm{def:m2:government-bonds:duration}{DV01}, et nommer la position sur la courbe que crée
la couverture.}
\end{solution}

\begin{solution}{iq:m2:government-bonds:5}
Échéances en fin de mois (règle de fin de mois, y compris février lors des
années bissextiles)~; règlement à une date de \omterm{def:m2:government-bonds:coupon}{coupon} (couru nul)~; règlement pendant une
période ex-dividende (couru négatif, pour les gilts)~; la \omterm{def:m2:government-bonds:daycount}{convention de décompte des jours} de chaque marché
(actual/actual pour les Treasuries et les gilts, 30/360 lorsque les conditions d'une obligation
le prévoient)~;
\omterm{def:m2:government-bonds:coupon}{coupons} annuels contre semestriels~; premiers \omterm{def:m2:government-bonds:coupon}{coupons} courts et longs~; les
exemples corrigés publiés par le régulateur et par la bourse, reproduits au
chiffre près~; allers-retours entre prix et rendement.
\iqlookfor{exemples corrigés publiés utilisés comme oracles, et cas limites du
calendrier.}
\end{solution}

\begin{solution}{iq:m2:government-bonds:6}
Une perpétuité versant $c$ par an a pour prix $P = c/y$. Alors $dP/dy = -c/y^2$ et
la \omterm{def:m2:government-bonds:duration}{duration modifiée} vaut $-P'/P = 1/y = 25$ ans à 4~\%. Sa \omterm{def:m2:government-bonds:duration}{duration de Macaulay}
vaut $(1 + y)/y = 26$ ans~: le temps moyen des paiements pondéré par leur valeur actuelle,
$\sum_k k v^k / \sum_k v^k = 1/(1 - v) = (1+y)/y$.
\iqlookfor{la forme fermée et l'écart entre les deux durations.}
\end{solution}
